% Appendix G: Experiment 4 (natural context, usage test). Written 2026-09-11 from
% code/experiments/exp4_natural_context.py, exp4_train.py, results/exp4*/ and
% REFOCUS_PLAN_2026-09-10.md section 8 (pre-registration, amendments 8.3a, 8.4, 8.5).

\subsection{Purpose and registration}
Experiment~4 tests, with natural neighbourhoods, the inference that motivated this
work: that a spatial--spectral model's insensitivity to spectral compression shows the
spectral dimension to be largely redundant \citep{oleary2026spatial,mueller2023dimred}.
It is a usage test, not a mechanism test: no cue is constructed or calibrated. Its
predictions, arms and disposition rules were fixed before the first run and are
reproduced in the pre-registration file; two amendments were recorded before the
arms they concern were run (the class-weighted loss, and the no-bottleneck arms
added after the authors' code was checked). Every registered outcome is reported,
including the ones that failed.

\subsection{Data and natural patches}
The cores, patient-level fold-0 split and four classes are those of
Appendix~\ref{app:exp3}. Two copies of the spectra are used: the non-denoised copy
(full rank; the pre-registered primary set) and the companion pipeline's
PCA-23-denoised copy (the field-standard preprocessing, closer to the pipeline of
\citet{oleary2026spatial}; a replication with identical arms and predictions,
registered before its cache was built). A patch is a labelled centre pixel with its
eight recorded neighbours, all nine inside the tissue mask; neighbours keep their own
labels or lack of one. Per core and class, up to $150$ (training) or $100$
(validation, test) patches were sampled with a fixed seed: $19{,}650$ training patches
from $79$ cores, $3{,}700$ validation from $20$, $2{,}700$ test from $16$ ($54$ cores
carry no label of the four classes). The neighbourhoods are homogeneous: no patch
contains a neighbour of a different labelled class, $88\%$ of neighbours share the
centre's label and $12\%$ are unlabelled. Natural $3\times3$ context on this data is
therefore a redundant copy of the centre's class signal, not a separate cue.

\subsection{Stage A: per-pixel information against the number of components}
On the standardized centre spectra, with PCA fitted on training patients, three
per-pixel classifiers were fitted on training centres for each number of components
$k$ and scored on validation and test patients: a shrinkage Fisher discriminant,
balanced logistic regression and a $256$-unit ReLU MLP with a fixed $40$-epoch
schedule and three seeds, with no selection on the evaluation patients
(Tables~\ref{tab:exp4_pixel_nodenoise} and~\ref{tab:exp4_pixel_pca23}). The
registered gate was a gain of at least $0.02$ macro-F1 from $k=16$ to the full
spectrum on both evaluation sides. On the non-denoised copy the gains are
$0.065$--$0.107$ (LDA $0.560\to0.651$ validation, $0.538\to0.645$ test; logistic
$0.595\to0.673$, $0.571\to0.658$; MLP $0.640\to0.705$, $0.594\to0.677$). The
dependence on $k$ is not monotone: components $9$--$16$ lower test macro-F1 for all
three classifiers before later components recover it, and the best MLP value is at
$k=128$ ($0.717$ validation). Sixteen principal components are therefore not enough for
per-pixel classification of these spectra by these classifiers. The top $16$
components carry $81.7\%$ of the standardized variance ($23$: $86.5\%$; $64$:
$97.7\%$).

\subsection{Stage B: joint training on natural patches}
The model is that of Experiment~3 with four logits: linear encoder $942\to12$ (or
$k\to12$ in the compression arms, the projection applied to all nine pixels), ReLU
patch head of width $M=32$, class-weighted mean cross-entropy (weights proportional to
inverse class frequency, normalized to mean one), full-batch plain gradient descent
with one global rate $10^{-3}$, a budget of $40{,}000$ steps, paired initialization
within seed, three seeds. Arms: \emph{natural} context; \emph{natural, 16 PCs}
(inputs projected onto the top $16$ training-centre principal directions);
\emph{shuffled} (the eight neighbours replaced by random training centres drawn
independently of the centre label: the uninformative-context comparator, jointly
trained); \emph{shuffled, 16 PCs}; \emph{frozen} (random encoder, head trained on
natural patches); \emph{natural, $\kappa=1/256$} (whole-head rate multiplier); and,
under amendment 8.5, \emph{no bottleneck} (identity encoder; the head reads the nine
raw standardized spectra) with natural and with shuffled context, the analogue of the
``fixed'' models of \citet{oleary2026spatial}. Evaluation on validation patients (test
reported separately), each centre held fixed: \emph{iid} natural patches;
\emph{context-random} (neighbours replaced by random centres from other cores of the
same side, independent labels); \emph{homogeneous} (all nine pixels equal to the
centre). The probe is a four-class discriminant on the encoder output of centre
spectra, fitted on $4{,}000$ held-out training centres. Retraining follows the
protocol of Section~\ref{sec:synthetic}: a fresh paired-initialization head trained
for $20{,}000$ steps on context-random training patches from a frozen saved encoder,
scored under context-random evaluation. Because the frozen and 16-component arms
plateau above every registered matched-fit threshold (loss $0.61$ and $0.53$
respectively), the registered fallback applies and the primary comparison is at
budget end; matched values at loss $0.5$ are given in the tables.

\paragraph{Results, non-denoised copy (Tables~\ref{tab:exp4_val_nodenoise},
\ref{tab:exp4_test_nodenoise}, \ref{tab:exp4_recovery_nodenoise}).}
\emph{P7, insensitivity of the natural-context model to PCA-16 compression: not
met.} At budget end the natural-context model scores $0.732$ macro-F1 in distribution
with the full spectrum and $0.657$ with $16$ components on validation patients
($0.652$ and $0.611$ on test); the full-spectrum model keeps fitting (loss $0.19$)
where the compressed one plateaus ($0.42$). At matched loss $0.5$ the two are
indistinguishable ($0.643$ and $0.636$). \emph{P8, reliance on the neighbourhood:
met.} Under context-random evaluation the natural-context model falls to $0.234$
against $0.680$ for the shuffled comparator and $0.705$ for the per-pixel MLP.
\emph{P9, where the information stopped: head-level.} The retrained head recovers
$0.681$ from the natural-context encoder against $0.688$ from the comparator's, and the
two encoders' probes are equal ($0.672$ and $0.670$, from $0.52$ at initialization);
a random encoder yields only $0.494$ after retraining. With natural context the encoder
learns the class signal exactly as it does without, and the collapse under
context-random evaluation is the head's use of same-class neighbours. \emph{P10,
compression asymmetry: not established.} The comparator loses $0.111$ from PCA-16
($0.652\to0.541$), but so does the natural-context model, so the registered reading
(sensitive comparator, insensitive natural-context model) does not hold. \emph{The
secondary rate prediction is not met}: slowing the head to $\kappa=1/256$ lowers the
probe ($0.619$) and the recovered value ($0.644$), at budget end and at matched loss
$0.5$ ($0.616$ against $0.632$). With a context that is a redundant copy of the
spectral cue there is no competition for the head's rate to relieve.

\paragraph{Results, PCA-23-denoised copy (Tables~\ref{tab:exp4_pixel_pca23},
\ref{tab:exp4_val_pca23}, \ref{tab:exp4_test_pca23}, \ref{tab:exp4_recovery_pca23}).}
The replication registered before its cache was built agrees with the non-denoised copy
on every verdict. The per-pixel gate passes more strongly (MLP $0.665\to0.804$
validation, $0.634\to0.883$ test from $16$ to $942$ components). The natural-context
model is again sensitive to PCA-16 at budget end ($0.727$ against $0.683$ in every
seed; $0.839$ against $0.640$ on test) and indistinguishable at matched loss; it
collapses under context-random evaluation ($0.274$); its encoder's probe equals the
comparator's ($0.671$ and $0.676$ from $0.529$) and a retrained head recovers
$0.706$--$0.712$ against the comparator's $0.720$; the slow head does not raise the
probe consistently ($0.683$, two seeds up and one down). One clause fails on this copy:
the jointly trained comparator ($0.722$) stays $0.08$ below the per-pixel MLP ($0.804$),
as does the no-bottleneck comparator ($0.715$), so the gap is the head family and plain
gradient descent, not the bottleneck.

\paragraph{P11: the learned bottleneck against none.} The comparison
\citet{oleary2026spatial} actually made is a learned linear bottleneck against no
reduction, and it reproduces on both copies: with natural context the learned
twelve-dimensional bottleneck scores $0.732$ in distribution against $0.744$ without a
bottleneck (denoised copy $0.727$ against $0.744$; on its test patients the bottleneck
model is $0.037$ higher), and the shuffled-context comparators score $0.680$ against
$0.661$ under context-random evaluation ($0.722$ against $0.715$ denoised); the
no-bottleneck model fits faster, so the equality is not a fitting artefact. Taken with Stage A and P7, this separates two readings of ``a small set of
spectral features'': the class signal is low-dimensional as a learned linear subspace,
as any four-class discriminant is, and not as a small number of principal components;
the learned features draw on the whole spectrum, so compressibility is not redundancy.
Whether a spatial model leaves that information unused is a separate question, answered
here in the negative for natural $3\times3$ context (the model uses the full spectrum and
its encoder is not starved) and in the affirmative for the constructed context of
Section~\ref{sec:real}.

% generated by paper/figures_src/make_tables_exp4.py
\begin{table}[htbp]\centering\scriptsize
\caption{Experiment 4 (non-denoised copy), Stage A: four-class macro-F1 of per-pixel classifiers on the centre spectra against the number of principal components $k$ (PCA on standardized training-split centres). Shrinkage Fisher discriminant (LDA), balanced logistic regression (LR) and a 256-unit ReLU MLP with a fixed schedule (mean$\pm$sd over three seeds), on validation and test patients.}\label{tab:exp4_pixel_nodenoise}
\begin{tabular}{rcccccc}\toprule
$k$ & LDA (val) & LR (val) & MLP (val) & LDA (test) & LR (test) & MLP (test) \\ \midrule
2 & 0.477 & 0.482 & 0.468$\pm$0.023 & 0.479 & 0.485 & 0.465$\pm$0.024 \\
4 & 0.506 & 0.517 & 0.528$\pm$0.011 & 0.517 & 0.525 & 0.520$\pm$0.019 \\
8 & 0.558 & 0.571 & 0.616$\pm$0.011 & 0.596 & 0.606 & 0.621$\pm$0.007 \\
16 & 0.560 & 0.595 & 0.640$\pm$0.012 & 0.538 & 0.571 & 0.594$\pm$0.009 \\
23 & 0.592 & 0.615 & 0.654$\pm$0.004 & 0.563 & 0.587 & 0.612$\pm$0.012 \\
32 & 0.603 & 0.618 & 0.674$\pm$0.003 & 0.576 & 0.597 & 0.619$\pm$0.003 \\
64 & 0.635 & 0.656 & 0.696$\pm$0.005 & 0.607 & 0.627 & 0.643$\pm$0.002 \\
128 & 0.664 & 0.672 & 0.717$\pm$0.006 & 0.623 & 0.643 & 0.677$\pm$0.009 \\
256 & 0.650 & 0.668 & 0.703$\pm$0.003 & 0.629 & 0.655 & 0.680$\pm$0.011 \\
942 & 0.651 & 0.673 & 0.705$\pm$0.002 & 0.645 & 0.658 & 0.677$\pm$0.007 \\
\bottomrule\end{tabular}\end{table}
\begin{table}[htbp]\centering\scriptsize
\caption{Experiment 4 (non-denoised copy), Stage B, validation patients: mean$\pm$sd over seeds (number of seeds) at the matched threshold $L^\ast$ and at budget end (40{,}000 steps or training loss $0.10$). Probe = macro-F1 of a discriminant on the encoder output of centre spectra; gain relative to initialization. Macro-F1 of the trained model under iid (natural), context-random (neighbours replaced by pixels from other cores with independent labels) and homogeneous (all nine pixels equal to the centre) evaluation.}\label{tab:exp4_val_nodenoise}
\begin{tabular}{llrccccc}\toprule
arm & at & step & probe & probe gain & iid & ctx-random & homogeneous \\ \midrule
natural, no bottleneck & $L^\ast$ & 3842$\pm$321 & 0.645$\pm$0.010 & 0.000$\pm$0.000 & 0.719$\pm$0.002 & 0.293$\pm$0.008 & 0.689$\pm$0.002 (3) \\
natural & $L^\ast$ & 13171$\pm$676 & 0.666$\pm$0.006 & 0.144$\pm$0.026 & 0.723$\pm$0.005 & 0.279$\pm$0.008 & 0.681$\pm$0.003 (3) \\
shuffled, no bottleneck & $L^\ast$ & 12248$\pm$656 & 0.645$\pm$0.010 & 0.000$\pm$0.000 & 0.647$\pm$0.001 & 0.657$\pm$0.002 & 0.650$\pm$0.008 (3) \\
\midrule
natural, no bottleneck & end & 30390$\pm$1099 & 0.645$\pm$0.010 & 0.000$\pm$0.000 & 0.744$\pm$0.005 & 0.249$\pm$0.007 & 0.717$\pm$0.002 (3) \\
natural & end & 40000$\pm$0 & 0.672$\pm$0.006 & 0.150$\pm$0.025 & 0.732$\pm$0.009 & 0.234$\pm$0.003 & 0.705$\pm$0.003 (3) \\
natural, $\kappa=1/256$ & end & 40000$\pm$0 & 0.628$\pm$0.015 & 0.106$\pm$0.024 & 0.652$\pm$0.012 & 0.296$\pm$0.018 & 0.631$\pm$0.007 (3) \\
natural, 16 PCs & end & 40000$\pm$0 & 0.574$\pm$0.010 & 0.031$\pm$0.020 & 0.657$\pm$0.004 & 0.201$\pm$0.024 & 0.620$\pm$0.007 (3) \\
shuffled, no bottleneck & end & 40000$\pm$0 & 0.645$\pm$0.010 & 0.000$\pm$0.000 & 0.675$\pm$0.005 & 0.661$\pm$0.006 & 0.674$\pm$0.008 (3) \\
shuffled (comparator) & end & 40000$\pm$0 & 0.670$\pm$0.010 & 0.148$\pm$0.012 & 0.652$\pm$0.008 & 0.680$\pm$0.006 & 0.648$\pm$0.012 (3) \\
shuffled, 16 PCs & end & 40000$\pm$0 & 0.557$\pm$0.011 & 0.015$\pm$0.014 & 0.541$\pm$0.004 & 0.583$\pm$0.011 & 0.540$\pm$0.008 (3) \\
natural, frozen encoder & end & 40000$\pm$0 & 0.521$\pm$0.019 & 0.000$\pm$0.000 & 0.576$\pm$0.015 & 0.264$\pm$0.022 & 0.540$\pm$0.016 (3) \\
\bottomrule
\end{tabular}\end{table}
\begin{table}[htbp]\centering\scriptsize
\caption{Experiment 4 (non-denoised copy), Stage B, test patients: mean$\pm$sd over seeds (number of seeds) at the matched threshold $L^\ast$ and at budget end (40{,}000 steps or training loss $0.10$). Probe = macro-F1 of a discriminant on the encoder output of centre spectra; gain relative to initialization. Macro-F1 of the trained model under iid (natural), context-random (neighbours replaced by pixels from other cores with independent labels) and homogeneous (all nine pixels equal to the centre) evaluation.}\label{tab:exp4_test_nodenoise}
\begin{tabular}{llrccccc}\toprule
arm & at & step & probe & probe gain & iid & ctx-random & homogeneous \\ \midrule
natural, no bottleneck & $L^\ast$ & 3842$\pm$321 & 0.651$\pm$0.011 & 0.000$\pm$0.000 & 0.648$\pm$0.007 & 0.255$\pm$0.008 & 0.630$\pm$0.004 (3) \\
natural & $L^\ast$ & 13171$\pm$676 & 0.614$\pm$0.011 & 0.097$\pm$0.027 & 0.655$\pm$0.002 & 0.242$\pm$0.007 & 0.627$\pm$0.003 (3) \\
shuffled, no bottleneck & $L^\ast$ & 12248$\pm$656 & 0.651$\pm$0.011 & 0.000$\pm$0.000 & 0.604$\pm$0.018 & 0.615$\pm$0.004 & 0.601$\pm$0.014 (3) \\
\midrule
natural, no bottleneck & end & 30390$\pm$1099 & 0.651$\pm$0.011 & 0.000$\pm$0.000 & 0.667$\pm$0.002 & 0.220$\pm$0.013 & 0.630$\pm$0.002 (3) \\
natural & end & 40000$\pm$0 & 0.610$\pm$0.009 & 0.093$\pm$0.025 & 0.652$\pm$0.004 & 0.206$\pm$0.004 & 0.618$\pm$0.004 (3) \\
natural, $\kappa=1/256$ & end & 40000$\pm$0 & 0.601$\pm$0.009 & 0.084$\pm$0.018 & 0.604$\pm$0.003 & 0.292$\pm$0.024 & 0.580$\pm$0.006 (3) \\
natural, 16 PCs & end & 40000$\pm$0 & 0.548$\pm$0.007 & 0.025$\pm$0.006 & 0.611$\pm$0.005 & 0.190$\pm$0.016 & 0.562$\pm$0.004 (3) \\
shuffled, no bottleneck & end & 40000$\pm$0 & 0.651$\pm$0.011 & 0.000$\pm$0.000 & 0.616$\pm$0.016 & 0.616$\pm$0.007 & 0.616$\pm$0.012 (3) \\
shuffled (comparator) & end & 40000$\pm$0 & 0.616$\pm$0.006 & 0.099$\pm$0.011 & 0.597$\pm$0.006 & 0.628$\pm$0.004 & 0.594$\pm$0.014 (3) \\
shuffled, 16 PCs & end & 40000$\pm$0 & 0.535$\pm$0.000 & 0.013$\pm$0.012 & 0.518$\pm$0.009 & 0.557$\pm$0.004 & 0.517$\pm$0.009 (3) \\
natural, frozen encoder & end & 40000$\pm$0 & 0.517$\pm$0.016 & 0.000$\pm$0.000 & 0.568$\pm$0.014 & 0.235$\pm$0.004 & 0.518$\pm$0.022 (3) \\
\bottomrule
\end{tabular}\end{table}
\begin{table}[htbp]\centering\scriptsize
\caption{Experiment 4 (non-denoised copy), head retraining: a fresh paired-initialization head trained for 20{,}000 GD steps on context-random training patches from a frozen saved encoder; macro-F1 under context-random and iid evaluation on validation and test patients (mean$\pm$sd over seeds; n).}\label{tab:exp4_recovery_nodenoise}
\begin{tabular}{llcccc}\toprule
encoder from & saved & ctx-random (val) & iid (val) & ctx-random (test) & iid (test) \\ \midrule
natural & random init & 0.494$\pm$0.036 & 0.482$\pm$0.020 & 0.481$\pm$0.044 & 0.472$\pm$0.049 (3) \\
natural & at $L^\ast$ & 0.678$\pm$0.006 & 0.667$\pm$0.009 & 0.624$\pm$0.005 & 0.612$\pm$0.006 (3) \\
natural & at end & 0.681$\pm$0.010 & 0.673$\pm$0.020 & 0.617$\pm$0.008 & 0.600$\pm$0.004 (3) \\
natural, $\kappa=1/256$ & at end & 0.644$\pm$0.008 & 0.634$\pm$0.019 & 0.604$\pm$0.002 & 0.578$\pm$0.018 (3) \\
natural, 16 PCs & at end & 0.594$\pm$0.012 & 0.574$\pm$0.015 & 0.568$\pm$0.002 & 0.544$\pm$0.016 (3) \\
shuffled & at end & 0.688$\pm$0.007 & 0.660$\pm$0.011 & 0.631$\pm$0.007 & 0.607$\pm$0.012 (3) \\
\bottomrule\end{tabular}\end{table}
\begin{table}[htbp]\centering\scriptsize
\caption{Experiment 4 (PCA-23-denoised copy), Stage A: four-class macro-F1 of per-pixel classifiers on the centre spectra against the number of principal components $k$ (PCA on standardized training-split centres). Shrinkage Fisher discriminant (LDA), balanced logistic regression (LR) and a 256-unit ReLU MLP with a fixed schedule (mean$\pm$sd over three seeds), on validation and test patients.}\label{tab:exp4_pixel_pca23}
\begin{tabular}{rcccccc}\toprule
$k$ & LDA (val) & LR (val) & MLP (val) & LDA (test) & LR (test) & MLP (test) \\ \midrule
2 & 0.506 & 0.511 & 0.496$\pm$0.005 & 0.512 & 0.524 & 0.507$\pm$0.006 \\
4 & 0.533 & 0.538 & 0.537$\pm$0.014 & 0.545 & 0.540 & 0.538$\pm$0.013 \\
8 & 0.496 & 0.552 & 0.601$\pm$0.008 & 0.524 & 0.597 & 0.615$\pm$0.005 \\
16 & 0.574 & 0.598 & 0.665$\pm$0.005 & 0.543 & 0.586 & 0.634$\pm$0.007 \\
23 & 0.624 & 0.649 & 0.726$\pm$0.007 & 0.617 & 0.658 & 0.691$\pm$0.009 \\
32 & 0.683 & 0.710 & 0.773$\pm$0.005 & 0.667 & 0.722 & 0.739$\pm$0.001 \\
64 & 0.684 & 0.708 & 0.775$\pm$0.006 & 0.697 & 0.748 & 0.791$\pm$0.006 \\
128 & 0.764 & 0.765 & 0.800$\pm$0.002 & 0.836 & 0.852 & 0.850$\pm$0.008 \\
256 & 0.705 & 0.773 & 0.802$\pm$0.003 & 0.857 & 0.859 & 0.880$\pm$0.005 \\
942 & 0.695 & 0.773 & 0.804$\pm$0.002 & 0.807 & 0.858 & 0.883$\pm$0.005 \\
\bottomrule\end{tabular}\end{table}
\begin{table}[htbp]\centering\scriptsize
\caption{Experiment 4 (PCA-23-denoised copy), Stage B, validation patients: mean$\pm$sd over seeds (number of seeds) at the matched threshold $L^\ast$ and at budget end (40{,}000 steps or training loss $0.10$). Probe = macro-F1 of a discriminant on the encoder output of centre spectra; gain relative to initialization. Macro-F1 of the trained model under iid (natural), context-random (neighbours replaced by pixels from other cores with independent labels) and homogeneous (all nine pixels equal to the centre) evaluation.}\label{tab:exp4_val_pca23}
\begin{tabular}{llrccccc}\toprule
arm & at & step & probe & probe gain & iid & ctx-random & homogeneous \\ \midrule
natural, no bottleneck & $L^\ast$ & 1380$\pm$140 & 0.756$\pm$0.008 & 0.000$\pm$0.000 & 0.714$\pm$0.006 & 0.302$\pm$0.012 & 0.705$\pm$0.006 (3) \\
natural & $L^\ast$ & 4141$\pm$41 & 0.676$\pm$0.017 & 0.147$\pm$0.006 & 0.725$\pm$0.002 & 0.303$\pm$0.014 & 0.699$\pm$0.005 (3) \\
shuffled, no bottleneck & $L^\ast$ & 6522$\pm$247 & 0.756$\pm$0.008 & 0.000$\pm$0.000 & 0.679$\pm$0.008 & 0.698$\pm$0.003 & 0.676$\pm$0.009 (3) \\
shuffled (comparator) & $L^\ast$ & 16013$\pm$1176 & 0.679$\pm$0.008 & 0.150$\pm$0.009 & 0.689$\pm$0.011 & 0.720$\pm$0.005 & 0.688$\pm$0.019 (3) \\
\midrule
natural, no bottleneck & end & 10371$\pm$474 & 0.756$\pm$0.008 & 0.000$\pm$0.000 & 0.744$\pm$0.001 & 0.280$\pm$0.010 & 0.730$\pm$0.003 (3) \\
natural & end & 20853$\pm$756 & 0.671$\pm$0.013 & 0.142$\pm$0.002 & 0.727$\pm$0.004 & 0.274$\pm$0.018 & 0.712$\pm$0.004 (3) \\
natural, $\kappa=1/256$ & end & 40000$\pm$0 & 0.683$\pm$0.016 & 0.155$\pm$0.015 & 0.697$\pm$0.010 & 0.299$\pm$0.018 & 0.685$\pm$0.013 (3) \\
natural, 16 PCs & end & 40000$\pm$0 & 0.582$\pm$0.002 & 0.024$\pm$0.024 & 0.683$\pm$0.006 & 0.187$\pm$0.014 & 0.637$\pm$0.005 (3) \\
shuffled, no bottleneck & end & 34031$\pm$890 & 0.756$\pm$0.008 & 0.000$\pm$0.000 & 0.717$\pm$0.018 & 0.715$\pm$0.004 & 0.715$\pm$0.017 (3) \\
shuffled (comparator) & end & 40000$\pm$0 & 0.676$\pm$0.002 & 0.147$\pm$0.009 & 0.703$\pm$0.003 & 0.722$\pm$0.004 & 0.700$\pm$0.004 (3) \\
shuffled, 16 PCs & end & 40000$\pm$0 & 0.569$\pm$0.000 & 0.011$\pm$0.021 & 0.564$\pm$0.024 & 0.588$\pm$0.007 & 0.556$\pm$0.023 (3) \\
natural, frozen encoder & end & 40000$\pm$0 & 0.529$\pm$0.011 & 0.000$\pm$0.000 & 0.594$\pm$0.026 & 0.261$\pm$0.016 & 0.556$\pm$0.015 (3) \\
\bottomrule
\end{tabular}\end{table}
\begin{table}[htbp]\centering\scriptsize
\caption{Experiment 4 (PCA-23-denoised copy), Stage B, test patients: mean$\pm$sd over seeds (number of seeds) at the matched threshold $L^\ast$ and at budget end (40{,}000 steps or training loss $0.10$). Probe = macro-F1 of a discriminant on the encoder output of centre spectra; gain relative to initialization. Macro-F1 of the trained model under iid (natural), context-random (neighbours replaced by pixels from other cores with independent labels) and homogeneous (all nine pixels equal to the centre) evaluation.}\label{tab:exp4_test_pca23}
\begin{tabular}{llrccccc}\toprule
arm & at & step & probe & probe gain & iid & ctx-random & homogeneous \\ \midrule
natural, no bottleneck & $L^\ast$ & 1380$\pm$140 & 0.889$\pm$0.025 & 0.000$\pm$0.000 & 0.705$\pm$0.008 & 0.274$\pm$0.007 & 0.691$\pm$0.005 (3) \\
natural & $L^\ast$ & 4141$\pm$41 & 0.673$\pm$0.009 & 0.105$\pm$0.026 & 0.702$\pm$0.006 & 0.271$\pm$0.015 & 0.688$\pm$0.011 (3) \\
shuffled, no bottleneck & $L^\ast$ & 6522$\pm$247 & 0.889$\pm$0.025 & 0.000$\pm$0.000 & 0.689$\pm$0.002 & 0.694$\pm$0.005 & 0.680$\pm$0.008 (3) \\
shuffled (comparator) & $L^\ast$ & 16013$\pm$1176 & 0.693$\pm$0.006 & 0.125$\pm$0.024 & 0.698$\pm$0.016 & 0.720$\pm$0.004 & 0.693$\pm$0.022 (3) \\
\midrule
natural, no bottleneck & end & 10371$\pm$474 & 0.889$\pm$0.025 & 0.000$\pm$0.000 & 0.802$\pm$0.005 & 0.248$\pm$0.008 & 0.777$\pm$0.001 (3) \\
natural & end & 20853$\pm$756 & 0.737$\pm$0.013 & 0.169$\pm$0.029 & 0.839$\pm$0.005 & 0.226$\pm$0.013 & 0.793$\pm$0.008 (3) \\
natural, $\kappa=1/256$ & end & 40000$\pm$0 & 0.674$\pm$0.013 & 0.105$\pm$0.008 & 0.678$\pm$0.006 & 0.293$\pm$0.027 & 0.669$\pm$0.010 (3) \\
natural, 16 PCs & end & 40000$\pm$0 & 0.553$\pm$0.002 & 0.007$\pm$0.031 & 0.640$\pm$0.003 & 0.177$\pm$0.006 & 0.582$\pm$0.003 (3) \\
shuffled, no bottleneck & end & 34031$\pm$890 & 0.889$\pm$0.025 & 0.000$\pm$0.000 & 0.755$\pm$0.018 & 0.755$\pm$0.005 & 0.748$\pm$0.018 (3) \\
shuffled (comparator) & end & 40000$\pm$0 & 0.726$\pm$0.002 & 0.158$\pm$0.020 & 0.763$\pm$0.014 & 0.782$\pm$0.003 & 0.762$\pm$0.013 (3) \\
shuffled, 16 PCs & end & 40000$\pm$0 & 0.546$\pm$0.002 & 0.000$\pm$0.031 & 0.556$\pm$0.014 & 0.580$\pm$0.005 & 0.549$\pm$0.004 (3) \\
natural, frozen encoder & end & 40000$\pm$0 & 0.568$\pm$0.020 & 0.000$\pm$0.000 & 0.638$\pm$0.034 & 0.234$\pm$0.018 & 0.577$\pm$0.031 (3) \\
\bottomrule
\end{tabular}\end{table}
\begin{table}[htbp]\centering\scriptsize
\caption{Experiment 4 (PCA-23-denoised copy), head retraining: a fresh paired-initialization head trained for 20{,}000 GD steps on context-random training patches from a frozen saved encoder; macro-F1 under context-random and iid evaluation on validation and test patients (mean$\pm$sd over seeds; n).}\label{tab:exp4_recovery_pca23}
\begin{tabular}{llcccc}\toprule
encoder from & saved & ctx-random (val) & iid (val) & ctx-random (test) & iid (test) \\ \midrule
natural & random init & 0.503$\pm$0.018 & 0.492$\pm$0.019 & 0.523$\pm$0.053 & 0.508$\pm$0.072 (3) \\
natural & at $L^\ast$ & 0.712$\pm$0.004 & 0.703$\pm$0.006 & 0.705$\pm$0.005 & 0.689$\pm$0.017 (3) \\
natural & at end & 0.706$\pm$0.003 & 0.701$\pm$0.001 & 0.791$\pm$0.005 & 0.780$\pm$0.019 (3) \\
natural, $\kappa=1/256$ & at end & 0.702$\pm$0.010 & 0.705$\pm$0.010 & 0.698$\pm$0.010 & 0.679$\pm$0.012 (3) \\
natural, 16 PCs & at end & 0.595$\pm$0.011 & 0.592$\pm$0.015 & 0.584$\pm$0.002 & 0.573$\pm$0.015 (3) \\
shuffled & at $L^\ast$ & 0.726$\pm$0.001 & 0.703$\pm$0.012 & 0.732$\pm$0.008 & 0.716$\pm$0.013 (3) \\
shuffled & at end & 0.720$\pm$0.005 & 0.704$\pm$0.011 & 0.782$\pm$0.007 & 0.761$\pm$0.003 (3) \\
\bottomrule\end{tabular}\end{table}

